\documentclass[11pt,a4paper]{article}
\usepackage[utf8]{inputenc}
\usepackage[T1]{fontenc}
\usepackage{lmodern}
\usepackage[french]{babel}
\usepackage{amsmath,amssymb,mathtools}
\usepackage{icomma}
\usepackage[version=4]{mhchem}
\usepackage{siunitx}
\sisetup{output-decimal-marker={,},per-mode=symbol}
\usepackage{xcolor}
\usepackage{colortbl}
\usepackage{tikz}
\usepackage[most]{tcolorbox}
\usepackage{enumitem}
\usepackage{array}
\usepackage{booktabs}
\usepackage{fancyhdr}
\usepackage[hidelinks]{hyperref}
\usetikzlibrary{arrows.meta,positioning,calc}
\usepackage[a4paper,margin=2.1cm,headheight=26pt,headsep=10pt]{geometry}

\definecolor{primary}{RGB}{146,95,161}
\definecolor{primarydark}{RGB}{92,55,108}
\definecolor{acidecol}{RGB}{205,60,45}
\definecolor{basecol}{RGB}{40,110,200}
\definecolor{methcol}{RGB}{90,90,100}

\tcbset{boxbase/.style={enhanced,breakable,boxrule=0.9pt,arc=2.5pt,left=3mm,right=3mm,top=2.3mm,bottom=2.3mm,fonttitle=\bfseries,coltitle=white,toptitle=0.6mm,bottomtitle=0.6mm}}
\newtcolorbox[auto counter]{exobox}[1][]{boxbase,colback=primary!4,colframe=primary,title={\textsc{Exercice}~\thetcbcounter\ifx\relax#1\relax\relax\else\textmd{ --- \itshape #1}\fi}}
\newtcolorbox{consigne}{boxbase,colback=methcol!7,colframe=methcol,sharp corners}

\pagestyle{fancy}
\fancyhf{}
\fancyhead[L]{\small\textcolor{primarydark}{\textbf{Fiche 10} — Thème 1 : couples acide/base}}
\fancyhead[R]{\small\textcolor{primarydark}{\textsc{Exercices — Moyen}}}
\fancyfoot[C]{\small\thepage}
\renewcommand{\headrulewidth}{0.8pt}
\renewcommand{\headrule}{\hbox to\headwidth{\color{primary}\leaders\hrule height \headrulewidth\hfill}}

\setlength{\parindent}{0pt}
\setlength{\parskip}{3pt}

\begin{document}

\begin{center}
{\LARGE\bfseries\color{primarydark} Thème 1 --- Niveau Moyen}\\[2pt]
{\color{primary}\itshape Identifier les deux couples d'une réaction, repérer les ampholytes}
\end{center}

\begin{consigne}
\textbf{Objectif du niveau.} Une réaction acide-base met en jeu \emph{deux} couples. Il faut savoir les
extraire, distinguer un vrai couple conjugué d'une fausse paire, et reconnaître une espèce qui peut jouer les
deux rôles (ampholyte).
\end{consigne}

\begin{exobox}[extraire les deux couples d'une réaction]
Pour chaque réaction, identifier les \textbf{deux couples acide/base conjugués} mis en jeu (préciser à chaque
fois qui est l'acide et qui est la base conjuguée de chaque couple).
\begin{enumerate}[label=\alph*),leftmargin=1.6em,itemsep=2pt]
  \item $\ce{CH3COOH + NH3 <=> CH3COO- + NH4+}$
  \item $\ce{HCO3- + OH- <=> CO3^{2-} + H2O}$
  \item $\ce{H2PO4- + NH3 <=> HPO4^{2-} + NH4+}$
\end{enumerate}
\end{exobox}

\begin{exobox}[la fausse paire]
Parmi les paires ci-dessous, l'une n'est \textbf{pas} un couple acide/base conjugués. La repérer et
\textbf{justifier} : rappeler qu'un couple conjugué se déduit par l'échange d'un \emph{seul} \ce{H+}.
\begin{enumerate}[label=\alph*),leftmargin=1.6em,itemsep=1pt]
  \item \ce{HClO}\,/\,\ce{ClO-}
  \item \ce{H3O+}\,/\,\ce{OH-}
  \item \ce{H2PO4-}\,/\,\ce{HPO4^{2-}}
  \item \ce{NH4+}\,/\,\ce{NH3}
\end{enumerate}
\end{exobox}

\begin{exobox}[montrer qu'une espèce est un ampholyte]
On considère l'ion hydrogénocarbonate \ce{HCO3-} (présent dans le sang).
\begin{enumerate}[label=\alph*),leftmargin=1.6em,itemsep=2pt]
  \item Écrire la réaction où \ce{HCO3-} joue le rôle d'\textbf{acide} vis-à-vis de l'eau, et donner le couple
        correspondant.
  \item Écrire la réaction où \ce{HCO3-} joue le rôle de \textbf{base} vis-à-vis de l'ion \ce{H3O+}, et donner
        le couple correspondant.
  \item Conclure : pourquoi dit-on que \ce{HCO3-} est un ampholyte ?
\end{enumerate}
\end{exobox}

\begin{exobox}[les bornes de l'eau]
Justifier, en essayant de construire le partenaire manquant, chacune des affirmations suivantes :
\begin{enumerate}[label=\alph*),leftmargin=1.6em,itemsep=2pt]
  \item « \ce{H3O+} n'a pas d'acide conjugué observable en solution aqueuse. »
  \item « L'eau \ce{H2O} est un ampholyte. » (donner ses deux couples)
\end{enumerate}
\end{exobox}

\begin{exobox}[repérer les ampholytes d'une famille]
On considère les quatre espèces issues des déprotonations successives de l'acide phosphorique :
\[ \ce{H3PO4},\quad \ce{H2PO4-},\quad \ce{HPO4^{2-}},\quad \ce{PO4^{3-}}. \]
\begin{enumerate}[label=\alph*),leftmargin=1.6em,itemsep=2pt]
  \item Parmi ces quatre espèces, lesquelles peuvent à la fois céder \emph{et} capter un \ce{H+} ? Ce sont les
        ampholytes.
  \item Pour l'espèce \ce{H2PO4-}, écrire les \textbf{deux couples} auxquels elle appartient.
\end{enumerate}
\end{exobox}

\end{document}
